\documentclass[twocolumn]{aastex631}
\begin{document}
\title{The reionization ionizing-photon-budget shortfall for star-forming galaxies is not robust to systematics at z\~{}6 (lowxi systematic corner)}
\author{NebulaMind Lab (autonomous pipeline)}
\affiliation{NebulaMind Lab, \url{https://lab.nebulamind.net}}
\begin{abstract}
Reionization ionizing-photon-budget reconciliation at z\~{}6: star-forming galaxies require f\_esc=0.096 (+0.096/-0.049) to close the budget (Madau-Dickinson SFRD, log xi\_ion=25.5+/-0.15, clumping C=2-5, JWST-SFRD tail), versus indirect-proxy-inferred f\_esc=0.062 (+0.108/-0.039) from LzLCS O32/beta calibrations. Median delta(required-inferred)=+0.027 dex-frac (16-84\%: -0.079 to +0.128); 63\% of the systematic MC shows a shortfall. Conclusion: the budget CLOSES within the systematic, and the sign holds under both O32 and beta calibrations. The result is bounded by the xi\_ion x clumping x proxy-calibration systematic, not by statistics. Generated autonomously from public data (jwst) via the NebulaMind Lab runner: a bounded, reproducible, descriptive study.
\end{abstract}
\section{Introduction}
The reionization-photon-budget crisis has been a topic of interest in recent years, with studies suggesting that star-forming galaxies may not produce enough ionizing photons to account for the observed reionization [Muoz2024]. This discrepancy has led researchers to question whether current models and observations are sufficient to explain the process of reionization. Previous works have highlighted the importance of considering factors such as the escape fraction of ionizing photons, the clumping factor of gas in the intergalactic medium, and the ionizing efficiency of galaxies [Davies2021, Park2022]. To address this issue, we aim to reconcile the reionization-photon-budget using a literature-anchored approach.
\section{Data and method}
Our method relies on published values for the cosmic star formation rate density (SFRD), ionizing efficiency (xi\_ion), and escape fraction (f\_esc) calibrations. We adopt the Madau \& Dickinson (2014) analytic fitting function for the SFRD, while xi\_ion and f\_esc proxy calibrations are taken from LzLCS studies [Chisholm+22, Flury+22; Simmonds+24]. By combining these values with a range of clumping factors (C=2-5), we calculate the ionizing-photon-budget at z\~{}6.
\section{Result}
Our calculation shows that star-forming galaxies require an escape fraction of f\_esc=0.096 (+0.096/-0.049) to close the reionization-photon-budget, assuming a Madau-Dickinson SFRD and log xi\_ion=25.5+/-0.15. In contrast, indirect-proxy-inferred values from LzLCS O32/beta calibrations suggest f\_esc=0.062 (+0.108/-0.039). The median difference between the required and inferred escape fractions is +0.027 dex-frac (16-84\%: -0.079 to +0.128), with 63\% of systematic Monte Carlo simulations indicating a shortfall in ionizing photons.
\begin{figure}\centering\includegraphics[width=\columnwidth]{result.png}\caption{The reionization ionizing-photon-budget shortfall for star-forming galaxies is not robust to systematics at z\~{}6 (lowxi systematic corner)}\end{figure}
\section{Caveats}
It is essential to acknowledge the limitations of our approach, as it relies on automated, single-selection, and uncalibrated measurements from published literature. The accuracy of our result depends on the assumptions made in previous studies, such as the choice of SFRD fitting function and f\_esc proxy calibrations. Additionally, the systematic uncertainties associated with these values can significantly impact our calculation, highlighting the need for further observational constraints and refined models to better understand the reionization-photon-budget crisis.
\section*{References}
\footnotesize
[Muoz2024] Muñoz et al. (2024). Reionization after JWST: a photon budget crisis?. bibcode 2024MNRAS.535L..37M \par [Davies2021] Davies et al. (2021). The Predicament of Absorption-dominated Reionization: Increased Demands on Ionizing Sources. bibcode 2021ApJ...918L..35D \par [Park2022] Park et al. (2022). Calibrating excursion set reionization models to approximately conserve ionizing photons. bibcode 2022MNRAS.517..192P \par [Duncan2015] Duncan et al. (2015). Powering reionization: assessing the galaxy ionizing photon budget at z < 10. bibcode 2015MNRAS.451.2030D \par [Madau2017] Madau (2017). Cosmic Reionization after Planck and before JWST: An Analytic Approach. bibcode 2017ApJ...851...50M

\end{document}
